\section{Joint compatibility in a hidden wiring family}

In a blind exchange, each laboratory prepares a pair of ports and keeps one.
An independent referee redistributes the sent ports by a hidden permutation.
Questions arrive after this exchange. The rewiring experiment of
Moradi Kalarde et al.~\cite{kalarde2026fermions} uses this information constraint
to compare information carriers. Here we consider a finite classical fragment
with independently prepared pairs and visible sender identities.

Write $r_i$ for the sender whose transmitted port arrives at receiver $i$.
For each allowed route $r$ a target specifies the complete set of conditional
probabilities $P_r(a\mid x)$. A witness chooses one joint probability law
$q_j(k_j,s_j)$ for each source $j$, with independent draws across sources.
At receiver $i$, the response is a stochastic kernel $D_{i,r_i}$ indexed by
its question, its retained value, and the incoming value. Source laws and
response tables are shared across routes. The compatibility equation is
\[
 P_r(a\mid x)=
 \sum_{k,s}\left(\prod_j q_j(k_j,s_j)\right)
             \prod_i D_{i,r_i}(a_i\mid x_i,k_i,s_{r_i}) .
\]
The quantifier over the witness precedes the quantifier over the route.
When $r_i=i$, the two observed ports belong to the same draw from $q_i$.

\begin{proposition}
If $r_i=r'_i$, any joint realization gives the same unconditional output
marginal at party $i$ for a fixed local question in routes $r,r'$.
\end{proposition}
\begin{proof}
Sum the compatibility equation over the other answers. Each other response
kernel sums to one for each latent assignment. The remaining expectation
depends only on $D_{i,r_i}$ and the joint law of $(k_i,s_{r_i})$. For
$r_i\ne i$ that law is the product of the two corresponding marginals of the
local preparations. For $r_i=i$ it is $q_i$ itself. Both the law and the
kernel agree across the two routes.
\end{proof}

For example, take routes $(2,0,1)$ and $(2,1,0)$ for receivers in order A,B,C.
If Alice must always output zero on the first route and one on the second at
the same question, each branch separately admits a deterministic witness.
Alice receives Charlie's port in both, so the proposition excludes a common
witness. For a compatible control let each party keep and send a copy of its
own fair bit, then output the XOR of its two observed bits. A three-cycle
yields a uniform even-parity triple. On a two-cycle the excluded party outputs
zero, and the two participating parties have the same fair output.

The evaluator verifies a supplied finite rational witness by direct summation.
The marginal rule can establish impossibility for arbitrary finite classical
port alphabets. Agreement of all tested marginals leaves the existence question
open in the absence of a matching witness. A claim about fermionic transmission
requires its own carrier algebra and admissibility rules.

\subsection{Six-mode fermionic reference}

For the three-party reference of Ref.~\cite{kalarde2026fermions}, treat the
five nonidentity permutations as routes and the four answers as
$0,1_+,1_-,2$. The initial state has one particle per origin, with equal
amplitudes on its kept and sent modes. Fix the output mode order to be
incoming then retained at each site, with the sites in alphabetical order.
Number sites and modes from zero. For a send choice $h_i\in\{0,1\}$ at origin
$i$, the occupied destination mode is
\[
 k_i(h)=\begin{cases}
  2i+1 & h_i=0,\\
  2r^{-1}(i) & h_i=1 .
 \end{cases}
\]
In the ordered occupation basis the amplitude acquires the sign
$s_h=(-1)^{\#\{(i,j):i<j,\ k_i(h)>k_j(h)\}}$ from anticommutation.
Thus each pair of choices $h,h'$ contributes $s_hs_{h'}/8$ to its density
matrix entry. The zero-phase ordinary tensor reference sets all $s_h=1$.

For each site use the one-particle basis $|01\rangle,|10\rangle$ and the
plus/minus outcome convention of Eqs.~(B.6), (B.17) and (B.18) of
Ref.~\cite{kalarde2026fermions}. The observable in this sector is
$(X+(-1)^x Z)/\sqrt{2}$ for Alice and $Z$ or $X$ for Bob and Charlie.
Vacuum and double occupation have distinct outcomes. All effects preserve
particle number, so their product in the final site ordering can be contracted
against the signed density entries above. The resulting probabilities are
in $\mathbb{Q}(\sqrt{2})$. For example, on route $ABC$, question $000$,
the probability of $(1_+,1_+,1_+)$ is $(2+\sqrt{2})/32$.

The reference calculation checks all $2560$ unconditioned probabilities
for positivity, normalization and nonsignalling in the late questions using
exact coefficient pairs. An independent numerical implementation propagates
the three occupied single-particle orbitals and computes squared determinants
for final detector configurations. The fermionic and ordinary tensor
references differ at $64$ entries in the full family. This discrepancy compares
two specified implementations. The paper's stronger exclusion theorem
quantifies over every admissible uniform qubit strategy and uses rigidity
across contexts.
